\usepackage{amsmath,amsfonts,amssymb,graphicx,xcolor,tikz,bm}
\usetikzlibrary{arrows.meta,shapes.geometric}
\usepackage[backend=biber]{biblatex}
\addbibresource{items/ALCO468-references.bib}
%\theoremstyle{plain}
\theoremstyle{definition}
\newtheorem{theorem}{Theorem}[section]
\newtheorem{proposition}[theorem]{Proposition}
\newtheorem{claim}[theorem]{Claim}
\newtheorem{lemma}[theorem]{Lemma}
\newtheorem{corollary}[theorem]{Corollary}
\newtheorem{problem}[theorem]{Problem}
\newtheorem{conjecture}[theorem]{Conjecture}
\newtheorem{definition}[theorem]{Definition}
\newtheorem{example}[theorem]{Example}
\newtheorem{algorithm}[theorem]{Algorithm} 
\newtheorem{procedure}[theorem]{Procedure}
\newtheorem{observation}[theorem]{Observation}
\newtheorem{case}{Case}
\newtheorem{xca}[theorem]{Exercise}
\newtheorem{consequence}[theorem]{Consequence}
\newtheorem{definition/theorem}[theorem]{Definition/Theorem}

\theoremstyle{remark}
\newtheorem{remark}[theorem]{Remark}
\newtheorem{notation}[theorem]{Notation}


\numberwithin{equation}{section}
\renewcommand{\labelitemi}{$\circ$} % A typesetting pun?!
\setcounter{MaxMatrixCols}{20}


%%%%%%blackboard letters
\newcommand{\bC}{\mathbb{C}}
\newcommand{\bN}{\mathbb{N}}
\newcommand{\bK}{\mathbb{K}}
\newcommand{\bP}{\mathbb{Z}_+} 
\newcommand{\bQ}{\mathbb{Q}}
\newcommand{\bR}{\mathbb{R}}
\newcommand{\bZ}{\mathbb{Z}}
\newcommand{\ka}{\kappa}
\newcommand{\la}{\lambda}
\newcommand{\bx}{\bm x}
\newcommand{\tre}{\textcolor{red}}
\newcommand{\tgr}{\textcolor{green}}
\newcommand{\dom}{\unrhd}
\newcommand{\domby}{\unlhd}

\DeclareMathOperator{\Asc}{Asc}
\DeclareMathOperator{\asc}{asc}
\DeclareMathOperator{\Inv}{Inv}
\DeclareMathOperator{\inv}{inv}
\DeclareMathOperator{\Fix}{Fix}
\DeclareMathOperator{\HT}{HT}
\DeclareMathOperator{\sgn}{sgn}

%:Sym definitions
\newcommand{\Sym}{\ensuremath{\operatorname{Sym}}}

%: compositions and tableaux
\newcommand\ol[1]{\overline{#1}}
\newcommand\ul[1]{\underline{#1}}
\newcommand\ch[1]{\check{#1}}
\newcommand\strongof[1]{{#1}^+}
\newcommand\partitionof[1]{\widetilde{#1}}
\newcommand\reverse[1]{{#1}^*}
\newcommand\ccomplement[1]{{#1}^c}
\newcommand\ctranspose[1]{{#1}^t}

% Others
\newcommand{\ra}{r_\alpha}
\newcommand{\rb}{r_\beta}
\newcommand{\bd}{P_{R,S}}
\newcommand{\lah}{\hat{\lambda}}
\newcommand{\al}{\alpha}


\newcommand*\circled[1]{\tikz[baseline=(char.base)]{
            \node[shape=circle,draw,inner sep=2pt] (char) {#1};}}
%%% tableau macros %%%
\newlength\cellsize \setlength\cellsize{15\unitlength}
\savebox2{%
\begin{picture}(15,15)
\put(0,0){\line(1,0){15}}
\put(0,0){\line(0,1){15}}
\put(15,0){\line(0,1){15}}
\put(0,15){\line(1,0){15}}
\end{picture}}
\newcommand\cellify[1]{\def\thearg{#1}\def\nothing{}%
\ifx\thearg\nothing
\vrule width0pt height\cellsize depth0pt\else
\hbox to 0pt{\usebox2\hss}\fi%
\vbox to 15\unitlength{
\vss
\hbox to 15\unitlength{\hss$#1$\hss}
\vss}}
\newcommand\tableau[1]{\vtop{\let\\=\cr
\setlength\baselineskip{-16000pt}
\setlength\lineskiplimit{16000pt}
\setlength\lineskip{0pt}
\halign{&\cellify{##}\cr#1\crcr}}}
\savebox3{%
\begin{picture}(15,15)
\put(0,0){\line(1,0){15}}
\put(0,0){\line(0,1){15}}
\put(15,0){\line(0,1){15}}
\put(0,15){\line(1,0){15}}
\end{picture}}
\newcommand\expath[1]{%
\hbox to 0pt{\usebox3\hss}%
\vbox to 15\unitlength{
\vss
\hbox to 15\unitlength{\hss$#1$\hss}
\vss}}
\newcommand\bas[1]{\omit \vbox to \cellsize{ \vss \hbox to \cellsize{\hss$#1$\hss} \vss}}

\newcommand{\emt}{\mbox{ }}
